\documentclass[11pt]{article}
\pagestyle{plain}

% \pgfplotsset{compat=1.18}
\usepackage[margin=1in]{geometry}
\usepackage{amsmath,amsfonts,amssymb,amsthm}
\usepackage[ruled,vlined,linesnumbered,commentsnumbered]{algorithm2e}
\SetKwInput{KwIn}{Input}
\SetKwInput{KwOut}{Output}
\SetKwInput{KwGlobal}{Global variable}
\SetKwInput{KwOracle}{Oracle}
\SetKwComment{Comment}{// }{}
\usepackage{graphicx}
\usepackage{enumerate}
\usepackage{bbm}
\usepackage{verbatim}
\usepackage{hyperref,color}
\usepackage[capitalize,nameinlink]{cleveref}
\usepackage[dvipsnames]{xcolor}
\hypersetup{
	colorlinks=true,
	pdfpagemode=UseNone,
	citecolor=OliveGreen,
	linkcolor=NavyBlue,
	urlcolor=Magenta,
	pdfstartview=FitW
}
\usepackage{appendix}
\usepackage{makecell}
\usepackage{tablefootnote}
\crefname{appsec}{Appendix}{Appendices}
\usepackage{tikz}
\usepackage{pgfplots}
\usepackage{xifthen}
\usepackage{subcaption}
\usepackage{hhline}
\usepackage{listings}

\theoremstyle{plain}
\newtheorem{theorem}{Theorem}[section]
\newtheorem{proposition}[theorem]{Proposition}
\newtheorem{lemma}[theorem]{Lemma}
\newtheorem{corollary}[theorem]{Corollary}
\newtheorem{conjecture}[theorem]{Conjecture}
\newtheorem{problem}[theorem]{Problem}
\newtheorem{claim}[theorem]{Claim}
\newtheorem{fact}[theorem]{Fact}

\theoremstyle{definition}
\newtheorem{definition}[theorem]{Definition}
\newtheorem{example}[theorem]{Example}
\newtheorem{observation}[theorem]{Observation}
\newtheorem{assumption}[theorem]{Assumption}
\newtheorem*{assumption*}{Assumption}
\newtheorem{condition}[theorem]{Condition}

\theoremstyle{remark}
\newtheorem{remark}[theorem]{Remark}

\crefname{lemma}{Lemma}{Lemmas}
\crefname{theorem}{Theorem}{Theorems}
\crefname{definition}{Definition}{Definitions}
\crefname{fact}{Fact}{Facts}
\crefname{claim}{Claim}{Claims}
\crefname{proposition}{Proposition}{Propositions}


\newcommand{\dif}{\,\mathrm{d}}
%\newcommand{\E}{\mathbb{E}}
%\newcommand{\Var}{\mathrm{Var}}
\newcommand{\Cov}{\mathrm{Cov}}
\newcommand{\MI}{\mathrm{I}}
%\newcommand{\Ent}{\mathrm{Ent}}
\newcommand{\ind}{\mathbbm{1}}
\newcommand{\allone}{\mathbf{1}}
\newcommand{\tr}{\mathrm{tr}}
\newcommand{\wt}{\mathrm{weight}}
\DeclareMathOperator*{\argmax}{arg\,max}
\newcommand{\norm}[1]{\left\lVert #1 \right\rVert}
\newcommand{\TV}[2]{\left\|#1 - #2\right\|_{\mathrm{TV}}}
\newcommand{\ceil}[1]{\left\lceil #1 \right\rceil}
\newcommand{\floor}[1]{\left\lfloor #1 \right\rfloor}
\newcommand{\ip}[2]{\left\langle #1 , #2 \right\rangle}
\newcommand{\grad}{\nabla}
\newcommand{\hessian}{\nabla^2}
\newcommand{\trans}{\intercal}
\newcommand{\diag}{\mathrm{diag}}
\newcommand{\poly}{\mathrm{poly}}
\newcommand{\dist}{\mathrm{dist}}
\newcommand{\sgn}{\mathrm{sgn}}
\newcommand{\fpras}{\mathsf{FPRAS}}
\newcommand{\fpaus}{\mathsf{FPAUS}}
\newcommand{\fptas}{\mathsf{FPTAS}}

\newcommand{\eps}{\varepsilon}

\newcommand{\N}{\mathbb{N}}
\newcommand{\Z}{\mathbb{Z}}
\newcommand{\Q}{\mathbb{Q}}
\newcommand{\R}{\mathbb{R}}
\newcommand{\C}{\mathbb{C}}

\renewcommand{\AA}{\mathcal{A}}
\newcommand{\BB}{\mathcal{B}}
\newcommand{\CC}{\mathcal{C}}
\newcommand{\DD}{\mathcal{D}}
\newcommand{\EE}{\mathcal{E}}
\newcommand{\FF}{\mathcal{F}}
\newcommand{\GG}{\mathcal{G}}
\newcommand{\HH}{\mathcal{H}}
\newcommand{\II}{\mathcal{I}}
\newcommand{\JJ}{\mathcal{J}}
\newcommand{\KK}{\mathcal{K}}
\newcommand{\LL}{\mathcal{L}}
\newcommand{\MM}{\mathcal{M}}
\newcommand{\NN}{\mathcal{N}}
\newcommand{\OO}{\mathcal{O}}
\newcommand{\PP}{\mathcal{P}}
\newcommand{\QQ}{\mathcal{Q}}
\newcommand{\RR}{\mathcal{R}}
\renewcommand{\SS}{\mathcal{S}}
\newcommand{\TT}{\mathcal{T}}
\newcommand{\UU}{\mathcal{U}}
\newcommand{\VV}{\mathcal{V}}
\newcommand{\WW}{\mathcal{W}}
\newcommand{\XX}{\mathcal{X}}
\newcommand{\YY}{\mathcal{Y}}
\newcommand{\ZZ}{\mathcal{Z}}

\newcommand{\KL}[2]{D_{\mathrm{KL}} \left( #1 \,\Vert\, #2 \right)}
\newcommand{\KLbig}[2]{D_{\mathrm{KL}} \left( #1 \,\big\Vert\, #2 \right)}
\newcommand{\KLBig}[2]{D_{\mathrm{KL}} \left( #1 \,\Big\Vert\, #2 \right)}

\newcommand{\SKL}[2]{D_{\mathrm{SKL}} \left( #1 \,\Vert\, #2 \right)}
\newcommand{\SKLbig}[2]{D_{\mathrm{SKL}} \left( #1 \,\big\Vert\, #2 \right)}
\newcommand{\SKLBig}[2]{D_{\mathrm{SKL}} \left( #1 \,\Big\Vert\, #2 \right)}

\newcommand{\qandq}{\quad\text{and}\quad}
\newcommand{\supp}{\mathrm{supp}}

\newcommand{\DKL}[2]{\-D_{\-{KL}}\left(#1 \parallel #2\right)}
\newcommand{\DTV}[2]{\-D_{\mathrm{TV}}\left({#1},{#2}\right)}
% \newcommand{\dist}{\mathrm{dist}}
\newcommand{\e}{\mathrm{e}}
\renewcommand{\epsilon}{\varepsilon}
\renewcommand{\emptyset}{\varnothing}
% \newcommand{\norm}[1]{\left\Vert#1\right\Vert}
\newcommand{\set}[1]{\left\{#1\right\}}
\newcommand{\tuple}[1]{\left(#1\right)} 
% \newcommand{\eps}{\varepsilon}
\newcommand{\inner}[2]{\left\langle #1,#2\right\rangle}
\newcommand{\tp}{\tuple}
\newcommand{\ol}{\overline}
\newcommand{\abs}[1]{\left\vert#1\right\vert}
\newcommand{\ctp}[1]{\left\lceil#1\right\rceil}
\newcommand{\ftp}[1]{\left\lfloor#1\right\rfloor}
\newcommand{\Ind}[1]{\-{Ind}_{#1}}
\newcommand{\todo}[1]{{\color{red} TODO: #1}}

\newcommand{\Cor}{\Psi^{\mathrm{sym}}}
\newcommand{\cor}{\Psi^{\-{cor}}}
\newcommand{\Inf}{\Psi}
% \newcommand{\diag}{\-{diag}}

\def\*#1{\boldsymbol{#1}} % Use \*A for \mathbf{A}
\def\+#1{\mathcal{#1}} % Use \+A for \mathcal{A}
\def\-#1{\mathrm{#1}} % Use \-A for \mathrm{A}
\def\=#1{\mathbb{#1}} % Use \^A for \mathbb{A}
\def\!#1{\mathfrak{#1}} % Use \!A for \mathfrak{A}

\newcommand*\diff{\mathop{}\!\mathrm{d}}

% probability
% \DeclareMathOperator*{\oPr}{\mathbf{Pr}}
\def\oPr{\mathop{\mathrm{Pr}}}
\renewcommand{\Pr}[2][]{ \ifthenelse{\isempty{#1}}
  {\oPr\left[#2\right]}
  {\oPr_{#1}\left[#2\right]} } % Use \Pr[a]{b} for \mathbf{Pr}_a[b], \Pr{b} for  \mathbf{Pr}[b]

% expectation: relies on the package "dsfont"
% \DeclareMathOperator*{\oE}{\mathbb{E}}
\def\oE{\mathop{\mathbb{E}}}
\newcommand{\E}[2][]{ \ifthenelse{\isempty{#1}}
  {\oE\left[#2\right]}
  {\oE_{#1}\left[#2\right]} }

% variance
%\DeclareMathOperator*{\oVar}{\mathbf{Var}}
\def\oVar{\mathrm{Var}}
\newcommand{\Var}[2][]{ \ifthenelse{\isempty{#1}}
  {\oVar\left[#2\right]}
  {\oVar_{#1}\left[#2\right]} }

% entropy
% \DeclareMathOperator*{\oEnt}{\mathbf{Ent}}
\def\oEnt{\mathrm{Ent}}
\newcommand{\Ent}[2][]{ \ifthenelse{\isempty{#1}}
  {\oEnt\left[#2\right]}
  {\oEnt_{#1}\left[#2\right]} }

\newcommand{\PhiEnt}[2][]{ \ifthenelse{\isempty{#1}}
  {\oEnt^\phi\left[#2\right]}
  {\oEnt^\phi_{#1}\left[#2\right]} }


\newenvironment{Exampleparagraph}[1][]
{
    \par 
    \noindent 
    \textbf{\color{blue} #1}
    \par 
    \vspace{0.5em} 
    \itshape 
}
{
    \par 
    \vspace{1em} 
}

%\input{local-macro}

\newcommand{\RED}[1]{\textcolor{red}{#1}}

\newcommand{\bern}[1]{\mathrm{Bern}\left(#1\right)}
\newcommand{\pll}{\parallel}



\title{CS 7545 F26, Homework 1}
\date{Due on September 17, 2026}
\author{{\textbf{Zejia Chen}} \\ \textbf{GTID: 904178159}}

\pgfplotsset{compat=1.18} 
\begin{document}

\maketitle

\section*{Exercise 1}
Fix $i\neq j$.
By independence of entries of $\vec{x}^{(i)}$ and $\vec{x}^{(j)}$, we have that $\ind[\vec{x}^{(i)}_k \neq \vec{x}^{(j)}_k] \sim \-{Bernoulli}\tp{\frac 12}$ independently.
Therefore, by Chernoff's bound,
$$
  \Pr{\norm{\vec{x}^{(i)} - \vec x^{(j)}} < \sqrt d}
  = \Pr{\sum_{k=1}^d \ind[\vec{x}^{(i)}_k \neq \vec{x}^{(i)}_k] <  d/4}
$$
$$
  = \Pr{\frac d2 - \sum_{k=1}^d \ind[\vec{x}^{(i)}_k \neq \vec{x}^{(i)}_k] \geq  d/4} \leq \exp(-d/8) .
$$
By the union bound,
$$
  \Pr{\forall i\neq j: \norm{\vec{x}^{(i)} - \vec x^{(j)}} \geq \sqrt d} \geq 1 - \binom n2\exp(-d/8) \geq 1 - \delta,
$$
since $n\leq \sqrt{2\delta} \exp(d/16)$.
\section*{Exercise 2}
\textbf{(1)} Since $f_{Z_i}(z) = f_{Z_i}(-z)$ for any $z\in \R$, $\Pr{X_i \geq z} = \Pr{X_i \leq -z}$ for $z\in \R$.
Hence $X_i$ has the same distribution as $-X_i$ and by independence, the empirical average $A$ has the same distribution as $-A$.
Therefore, $\E{X_i} = \frac 12(\E{X_i}+\E{-X_i})=0$ and $\Pr{|A|\geq \eps} = \Pr{A \geq \eps} + \Pr{-A \geq \eps} = 2\Pr{A\geq \eps}$.
\\\textbf{(2)} By Hoeffding's inequality, we have that
$$
  \Pr{A \geq \eps} =  \Pr{A - \E{A}\geq \eps} 
  \leq \exp\tp{-\frac{n\eps^2}{2M^2}}.
$$
\textbf{(3)} The following holds
\begin{align*}
  \E{X_i^2} &= 2\int_{0}^{M}t^2 \frac {\exp(-t)}2 dt
  =  \left.-(t^2+2t+2)e^{-t}\right|_{0}^{M}
  =2-(M^2+2M+2)e^{-M} .
\end{align*}
\textbf{(4)} By Bernstein’s Inequality, we have that 
$$
  \Pr{A \geq \eps} =  \Pr{A - \E{A}\geq \eps} 
  \leq \exp\tp{-\frac{n\eps^2/2}{\frac 1n\sum_{i=1}^n \E{X_i^2}+\frac 13M\eps}} \leq \exp\tp{-\frac{n\eps^2}{4+2M\eps/3}},
$$
where the last inequality follows from $\E{X_i^2} \leq 2$ for $i\in[n]$.
\\\textbf{(5)}
Since $\eps\leq M$ and $M\geq 2$, we have that
$$
  \exp\tp{-\frac{n\eps^2}{4+2M\eps/3}} \leq \exp\tp{-\frac{n\eps^2}{M^2+2M\eps/3}} \leq \exp\tp{-\frac{n\eps^2}{2M^2}} .
$$
Therefore, Bernstein’s bound from Part \textbf{(4)} is tighter than the Hoeffding bound from Part \textbf{(2)}.
\\\textbf{(6)} We have that
\begin{align*}
  \E{\exp(tX_i)} &= \Pr{\abs{Z_i} > M} + \int_0^M \frac {e^{tx}+e^{-tx}}2\exp(-x) dx
  \\&= e^{-M} + \frac 12\tp{\frac {1-e^{(t-1)M}}{1-t} + \frac{1-e^{-(t+1)M}}{t+1}}
  \\&\leq e^{-M} + \frac 12\tp{\frac {1}{1-t} + \frac{1}{t+1} - 2e^{-M}} = \frac 1{1-t^2} .
\end{align*}
The inequality follows from the hint with $x=e^{-M}$.
\\\textbf{(7)} 
The following holds for $\eps > 0$ and $t\in (0,1)$,
\begin{align*}
  \Pr{A\geq \eps} &= \Pr{\exp(tA) \geq \exp(t\eps)}
  \\&\leq \frac{\E{\exp(tA)}}{\exp(t\eps)} = \frac{\displaystyle \prod_{i=1}^n \E{\exp(tX_i/n)}}{\exp(t\eps)} &&(\text{Markov's ineq.})
  \\&\leq \frac{(1-(t/n)^2)^{-n}}{\exp(t\eps)} = \exp\tp{-n\tp{\ln (1-\frac{t^2}{n^2}) + \frac{t\eps}{n}}}. &&(\text{Part \textbf{(6)}}) 
\end{align*}
Plugging $t = n(\sqrt{1+\eps^2} - 1)/\eps$ into the above inequality, we complete the proof.
\\\textbf{(8)}
Set $s = \sqrt{1+\eps^2} - 1$ and it is sufficient to show that
$$
  f(\eps) = \ln \frac 2{s+2} + s \geq \frac{2s(s+2)}{8+s(s+2)} .
$$
Let $g(x) = \ln(2/(x+2))+x-\frac{2x(x+2)}{8+x(x+2)}$.
We have that for $x\geq 0$,
$$
  g'(x) = \frac{x^2(x+1)(x^2+4x+20)}{(x+2)(x^2+2x+8)^2} \geq 0.
$$
Therefore $g$ is nondecreasing on $x\geq 0$ and hence, $g(x) \geq g(0) = 0$.
Since 
$$
  \exp(-nf(\eps)) \leq \exp\tp{-\frac{n\eps^2}{4+\frac 12 \eps^2}}\leq \exp\tp{-\frac{n\eps^2}{4+2\eps^2/3}}\leq \exp\tp{-\frac{n\eps^2}{4+2M\eps/3}}, 
$$
the bound from Part \textbf{(7)} is tighter than Bernstein’s inequality.
\section*{Exercise 3}
\textbf{(1)} Fix the values of $S_1,\ldots,S_{i-1}$.
Let $Y := S_{i} - S_{i-1}\mid S_1,\ldots,S_{i-1}$ and $Z := (Y+c)/2c$.
We have that $\E{Z} = \frac 12$ and $Z\in [0,1]$ a.s.
Therefore,
$$
  \E{\exp(tY)} = e^{-ct}\E{\exp(2tcZ)} 
  \leq e^{-ct}\E{1 - Z + e^{2tc}Z} = \frac{e^{tc}+e^{-tc}}2,
$$
where the inequality follows from the convexity of $x\mapsto \exp(x)$.
Since
$$
  \frac{e^x + e^{-x}}2 = \sum_{i=0}^{\infty}\frac{x^{2i}}{(2i)!}
  \leq \sum_{i=0}^{\infty}\frac{x^{2i}}{2^ii!} = e^{x^2/2},
$$
$$
  \E{\exp(tY)} \leq  \frac{e^{tc}+e^{-tc}}2 \leq \exp(t^2c^2/2).
$$
\textbf{(2)} By Part \textbf{(1)}, we have that
$$
  \E{\exp(t(S_{i} - S_{i-1}))\mid S_1,\ldots,S_{i-1}}
  = \E{\exp(tS_{i})\mid S_1,\ldots,S_{i-1}}\exp(-tS_{i-1})\leq \exp\tp{\frac{t^2c^2}{2}} .
$$
Hence, $\E{\exp(tS_{i})\mid S_1,\ldots,S_{i-1}} \leq \exp(\frac{t^2c^2}{2}+tS_{i-1})$.
Taking expectation on both sides, we conclude that
\begin{align*}
  \E{\exp(tS_i)} &= \E{\E{\exp(tS_i)\mid S_1,\ldots,S_{i-1}}}
  \\&\leq \E{\exp\tp{\frac{t^2c^2}{2}+tS_{i-1}}}
  = \exp\tp{\frac{t^2c^2}{2}}\E{\exp(tS_{i-1})} .
\end{align*}
\textbf{(3)} By Part \textbf{(2)}, we have that
\begin{align}\label{eq:E3_1}
  \E{\exp(tS_n)} \leq \exp\tp{\frac{t^2c^2}{2}}\E{\exp(tS_{n-1})}
  \leq \cdots \leq \exp\tp{\frac{t^2c^2}{2}}^n\E{\exp(tS_{0})} 
  =\exp\tp{\frac{nt^2c^2}{2}} .
\end{align}
By Markov's inequality and \Cref{eq:E3_1}, for any $t > 0$, the following holds,
\begin{align}\label{eq:E3_2}
  \Pr{S_n\geq \eps n} = \Pr{\exp(tS_n) \geq \exp(t\eps n)}
  \leq \frac{\E{\exp(tS_n)}}{\exp(t\eps n)}
  \leq \exp\tp{\frac{nt^2c^2}{2} - t\eps n} .
\end{align}
Plugging $t = \frac{\eps}{c^2}$ into \Cref{eq:E3_2}, we conclude that
$$
  \Pr{S_n\geq \eps n} \leq \exp\tp{-\frac{n\eps^2}{2c^2}} .
$$
\section*{Exercise 4}
\textbf{(1)} Since the samples in $S$ are i.i.d., Bernstein’s inequality and the union bound give
\begin{align*}
  \Pr{\exists h\in H: L(h) - \hat L_S(h) \geq \eps}
  &\leq |H|\max_{h\in H}\exp\tp{-\frac{3n\eps^2}{6L(h)(1-L(h))+2\eps}}
  \\&\leq |H|\exp\tp{-\frac{3n\eps^2}{6p(1-p)+2\eps}},
\end{align*}
since $p \leq 1/2$.
\\\textbf{(2)}
Set 
$$
  \eps = \sqrt{\frac{2p(1-p)\ln(|H|/\delta)}{n}} +\frac{2}{3n}\ln(|H|/\delta).
$$
By Part \textbf{(1)}, it is sufficient to show that
\begin{align}\label{eq:E4_1}
  |H|\exp\tp{-\frac{3n\eps^2}{6p(1-p)+2\eps}} \leq \delta
  \iff 3n\eps^2 - 2\ln\tp{\frac{|H|}{\delta}}\eps - 6\ln\tp{\frac{|H|}{\delta}}p(1-p)\geq 0 .
\end{align}
By solving the quadratic inequality, \Cref{eq:E4_1} holds if 
\begin{align}\label{eq:E4_2}
  \eps \geq \frac{\ln(|H|/\delta)}{3n} + \sqrt{\frac {(\ln(|H|/\delta))^2}{9n^2} + \frac{2p(1-p)\ln(|H|/\delta)}{n}}.
\end{align}
\Cref{eq:E4_2} immediately follows from $\sqrt{a+b} \leq \sqrt{a}+\sqrt{b}$ for $a,b \geq 0$.
\section*{Exercise 5}
\textbf{(1)}
By minimality, $L(\hat h_S) - L(h^*)\geq 0$.
Then Part \textbf{(1)} follows from Markov's inequality.
\\\textbf{(2)}
Let $A$ be the event that $L(\hat h_S) - L(h^*) \geq \eps$.
Since $L(\hat h_S) - L(h^*) \leq L(\hat h_S) \leq 1$, we have that
\begin{align*}
  \E{L(\hat h_S) - L(h^*)} &= \E{L(\hat h_S) - L(h^*)\mid \overline A}\Pr{\overline A} + \E{L(\hat h_S) - L(h^*)\mid A}\Pr{A}
  \\&\leq \E{L(\hat h_S) - L(h^*)\mid \overline A} + \Pr{A}
  \leq \eps + \delta .
\end{align*}
\textbf{(3)}
Since $\E{X} = \int_0^\infty \Pr{X > x} dx$, we have that for any $\eps > 0$,
\begin{align}
  E &= \int_0^\eps \Pr{L(\hat h_S) - L(h^*) > x} dx + \int_\eps^\infty \Pr{L(\hat h_S) - L(h^*) > x} dx \nonumber
  \\&\leq \int_0^\eps 1 dx + \int_\eps^\infty 2|H|\exp(-nx^2/2) dx \nonumber
  \\&=\eps + \frac{2|H|}{\sqrt n}\int_{\sqrt n \eps}^\infty \exp(-t^2/2) dt \nonumber
  \\&\leq \eps + \frac{2|H|}{n\eps}\exp(-n\eps^2/2). \label{eq:E5_1}
\end{align}
Plugging $\eps = \sqrt{\frac{2\ln(2|H|)}{n}}$ into \Cref{eq:E5_1}, we conclude that
$$
  E \leq \sqrt{\frac{2\ln(2|H|)}{n}} + \sqrt{\frac 1{2n\ln (2|H|)}} = \sqrt{\frac{2\ln(2|H|)}n} \left(1+\frac1{2\ln(2|H|)}\right).
$$
\end{document}